\section{前言}

\subsection{研究背景}

随着自动化技术和智能制造系统的发展，多机械臂协同作业逐渐成为仓储搬运、工业分拣、柔性装配和物流输送等场景中的重要研究方向。与单机械臂系统相比，多机械臂系统具有更强的并行作业能力，能够在同一作业区域内同时完成多个搬运任务，从而缩短整体作业时间，提高系统运行效率。然而，多机械臂系统也带来了更加复杂的调度问题：不同机械臂之间可能存在空间资源冲突，部分任务需要两只机械臂协同完成，不同任务之间存在执行顺序差异，同时系统还需要在完成时间、能耗和资源利用率之间进行综合权衡。

在实际搬运场景中，待搬运物块通常具有不同的尺寸、重量和目标位置。较小或较轻的物块可以由单只机械臂独立搬运，而较大或较重的物块则需要两只机械臂协同完成。与此同时，不同类别的物块需要被搬运到不同的收集盒中，任务的空间位置、搬运距离和协同要求都会影响最终调度结果。因此，多机械臂搬运问题并不是简单的任务平均分配问题，而是一个同时涉及任务分配、执行排序、资源互斥、协同模式选择和目标优化的组合优化问题。

运筹学中的整数规划、调度优化、分支定界和启发式算法为该类问题提供了有效的建模与求解工具。通过将机械臂视为有限作业资源，将物块搬运过程视为一组具有时间和资源约束的任务，可以把实际搬运问题转化为可求解的数学优化模型。在此基础上，结合精确求解算法与自定义搜索策略，不仅可以获得较优调度方案，还可以分析不同算法在求解质量、计算时间和适用场景方面的差异。

\subsection{研究现状}

现有多机械臂协同作业研究大多可以分为两个方向。一类研究侧重于机器人学层面的运动规划、轨迹控制、避障控制和动力学协调，重点解决机械臂如何稳定、准确地完成运动的问题；另一类研究则侧重于系统调度和任务规划，重点关注在多个任务和多个执行资源之间如何进行分配和排序。本文关注的是后者，即从运筹学角度研究多机械臂搬运任务的调度优化问题，而不涉及机械臂底层关节轨迹、运动学逆解和动力学控制。

在调度优化方法方面，整数规划模型能够较完整地描述任务选择、资源占用和时间先后关系，适合用于构建标准化数学模型。CP-SAT 求解器等通用优化工具能够对中小规模问题给出较高质量的解，但在问题规模增大或约束复杂时，求解时间可能明显增加。隐枚举法和分支定界思想则适合处理具有离散决策结构的问题，通过逐步构造候选解并利用当前最优上界、可行性判断和下界估计提前剪枝，可以避免对所有组合方案进行完全枚举。启发式算法虽然通常不能保证全局最优，但具有实现简单、计算速度快和适合快速生成可行解的特点，在实际工程问题中也具有较强应用价值。

对于多机械臂搬运任务而言，仅比较调度完成时间并不足以全面评价方案优劣。若机械臂通过提高运行速度来缩短搬运时间，往往会导致能耗增加；若增加机械臂数量，虽然可能提升并行搬运能力，但也可能带来额外的启用成本和能耗代价。因此，在基础整数规划模型之外，引入非线性速度--能耗分析，有助于进一步说明调度方案背后的时间收益与能耗代价之间的关系。

\subsection{研究意义}

本文研究的意义主要体现在三个方面。

首先，从实际问题抽象角度看，本文将多机械臂搬运场景转化为运筹学中的调度优化问题，明确区分了机械臂底层控制问题与任务层调度问题。通过设置任务分配变量、时间变量和模式选择变量，可以较清晰地描述单臂搬运、双臂协同搬运、多类别收集盒以及机械臂资源互斥等实际约束，使问题具有明确的数学表达。

其次，从算法设计角度看，本文不仅使用 CP-SAT 求解器作为标准优化工具，还设计了融合剪枝的隐枚举法、启发式算法和热启动后的隐枚举法。不同方法之间既有精确求解与近似求解的对比，也有通用求解器与自定义算法的对比。通过比较各方法在不同任务规模和不同场景下的运行结果，可以分析剪枝策略、启发式规则和热启动机制对求解效率的影响。

最后，从结果分析角度看，本文不仅关注最终调度方案是否可行，还进一步比较二臂与三臂场景下的完成时间、能耗和综合表现，并通过非线性优化分析讨论速度调整带来的时间--能耗权衡。这样的分析能够说明多机械臂系统并不是机械臂数量越多越好，也不是运行速度越快越优，而是需要根据任务规模、物块类型、协同需求和能耗约束进行综合判断。

\subsection{本文主要内容}

本文的主要研究内容如下。

第一，建立多机械臂搬运调度优化模型。根据多类别物块、不同目标收集盒、单臂与双臂协同搬运模式等实际条件，构建包含任务分配、时间一致性、资源互斥和序列转移约束的数学模型，并给出相应目标函数。

第二，设计并实现多种求解方法。本文分别采用 CP-SAT 求解器、融合剪枝的隐枚举法、启发式算法以及热启动后的隐枚举法进行求解，并从算法思想、实现方式和适用特点等角度进行说明。

第三，开展算例分析与方法对比。通过多组不同任务规模和不同机械臂数量的算例，比较各方法得到的最大完成时间、能耗、运行时间和负载均衡效果，分析不同算法的优缺点。

第四，进行非线性速度--能耗优化分析。在基础调度结果上，引入速度倍率变量，建立速度与作业时间、能耗之间的非线性关系，并结合 KKT 条件对优化结果进行验证，进一步讨论二臂与三臂模式选择中的时间收益和能耗代价。